\section{Selection Bias and Methodological Limitations}

\subsection{Retrospective cohort selection}

The cohort was defined using December 2022 liquidity and then applied to
observations beginning in 2016. This is retrospective cohort selection:
training-period observations are retained only for companies that would become
highly liquid by January 2023. It differs from a year-by-year point-in-time
universe that would allow liquidity to enter and exit the sample as trading
activity changed.

The resulting bias is potentially severe for studies of high-beta technology
stocks and high-volume reversal. A fixed liquid cohort over-represents
companies whose later success made them liquid and under-represents companies
that were liquid earlier but declined, merged, or delisted before the screen.
The survivorship and delisting-bias mechanisms are well documented in the
financial data literature \citep{brown1992survivorship,shumway1997delisting}.
All selected histories currently ending in 2025 reinforce this caveat; the
observed no-trailing-gap result is not evidence that the cohort was free of
survivorship selection.

\subsection{Universe and provider limitations}

Only \AcquiredSecurities{} of \UniverseSecurities{} SEC-derived securities
have acquired histories. Yahoo Finance is not a complete historical listing
database. Failed, renamed, delisted, or reused tickers can be missing or mapped
ambiguously. The project therefore treats this as a verified partial universe,
not a complete January 2023 census.

Listing and delisting dates remain unresolved for most master rows. The report
distinguishes leading observed-history gaps from exchange holidays and
weekends, but it does not relabel every leading gap as pre-listing. Two raw
placeholder intervals are corroborated by first valid observations; that is
evidence, not a certified exchange listing date.

\subsection{Vendor and corporate-action limitations}

Restoring the historical split basis requires future split metadata. The rule
is economically sensible for the intended price-volume multiplication, but it
depends on the completeness of Yahoo's corporate-action records. The current
policy does not guarantee that every reported volume was normalized between
splits, nor does it validate every dividend adjustment. Four corporate-action
case records remain subject to interpretation.

\subsection{Scope exclusions}

This supplement does not calculate final returns, create training features,
construct a rolling liquidity filter, split the sample chronologically, fit an
ML model, or run a backtest. Any extension to a rolling, point-in-time design
must separate eligibility from outcome data and must supply verified listing
and delisting information where available.
